\title{FlashAttention-3: Fast and Accurate Attention with Asynchrony and Low-precision}

\begin{abstract}
Serving as the foundational layer in the prevalent Transformer architecture, attention mechanisms represent a primary computational constraint for large language models and tasks requiring long-context reasoning. While \textsc{FlashAttention} advanced GPU-based attention by reducing memory traffic, it has not fully harnessed advancements in recent hardware platforms—\textsc{FlashAttention-2} notably utilizes just 35\% of the H100 GPU's capabilities. In this work, we introduce three optimization strategies tailored for Hopper GPUs: (1) utilizing the asynchrony between Tensor Cores and TMA engines to achieve concurrent computation and data transfer via warp-specialization, (2) integrating block-level matrix multiplication and softmax execution, and (3) employing block quantization with incoherent computation enabled by FP8 low-precision features native to the hardware. Through these contributions, our approach—\textsc{FlashAttention-3}—delivers a 1.5-2.0$\times$ increase in speed on H100 GPUs, reaching up to 740 TFLOPs/s at FP16 precision (representing 75\% hardware utilization) and nearly 1.2 PFLOPs/s at FP8. Our results further confirm that FP8 \textsc{FlashAttention-3} reduces numerical error by a factor of 2.6$\times$ compared to standard FP8 attention implementations.
\end{abstract}

\section{Introduction}
\label{sec:intro}

Within the Transformer architecture~\citep{vaswani2017attention}, the main limitation arises from the attention mechanism, as calculating self-attention scores between queries and keys leads to computational costs that scale quadratically with sequence length. Expanding the attention mechanism to handle longer contexts could open the door to a range of novel functionalities—including the ability to model and reason over multiple lengthy documents~\citep{guo2021longt5,shaham2022scrolls,peng2023yarn} or files in expansive codebases~\citep{roziere2023code, li2023starcoder}, as well as support additional modalities such as high-resolution visual data~\citep{chen2022scaling}, audio signals~\citep{gulati2020conformer}, or video streams~\citep{ho2022video}. Other use-cases such as supporting prolonged user interaction histories~\citep{sun2019bert4rec} or enabling agents with extended planning capabilities~\citep{yao2022react} would also benefit from improved scalability. These needs have prompted a surge of research into accelerating attention for long-context settings—ranging from methods based on approximation~\citep{katharopoulos2020transformers,choromanski2020rethinking,
tay2020efficient}, software-level optimizations (\citep{rabe2021self,dao2022flashattention,kwon2023efficient}), to investigations into alternative model architectures~\citep{peng2023rwkv,sun2023retentive,gu2023mamba}.

Our study extends upon the foundation established by \citet{dao2022flashattention}, which focused on designing exact-attention algorithms that leverage detailed knowledge of GPU architecture and operational paradigms in their conceptual framework. In their seminal contribution, Dao et al. proposed \textsc{FlashAttention}, an innovative approach utilizing tiling techniques to parallelize attention computation. This method effectively removes the need for intermediate accesses to the relatively slow global memory by merging all attention-related operations within a single GPU kernel.

Subsequently, \citet{dao2023flashattention2} introduced \textsc{FlashAttention-2}, enhancing the algorithm’s parallelism by spanning the sequence length dimension and executing the core of the forward pass across segmented blocks of the key and value matrices. This modification led to superior workload balancing and improved GPU resource utilization. Yet, our analysis reveals that \textsc{FlashAttention-2} continues to exhibit suboptimal performance on contemporary GPUs, such as the Hopper H100, with utilization rates markedly lower than those achieved by state-of-the-art matrix multiplication (GEMM) kernels—approximately 35\% compared to 80-90\%. This gap can be partly attributed to differences in the implementation, notably the absence of Hopper-specific instructions, favoring legacy Ampere instructions when deploying on Tensor Cores. Recent advancements, exemplified by ThunkerKitten~\citep{spector2024thunder} and cuDNN 9~\citep{cudnn9}, underscore that integrating Hopper-tailored instructions and tile-based methodologies can significantly accelerate attention operations and streamline code complexity.

At a more foundational level, the algorithm underlying \textsc{FlashAttention-2} operates within a straightforward synchronous paradigm, omitting any explicit integration of asynchronous execution or reduced-precision arithmetic. Asynchrony emerges as a consequence of specialized hardware intended to optimize critical components of machine learning workloads: matrix multiplication is accelerated by dedicated units such as Tensor Cores, while memory operations are managed by separate hardware like the Tensor Memory Accelerator (TMA), both functioning independently from general-purpose CUDA cores tasked with logic, integer, and floating-point processing. The ongoing shift toward low-precision formats, including FP8 on Hopper and FP4 on Blackwell, follows earlier adoption of FP16 on Pascal (2017) and BF16 on Ampere (2020). These approaches consistently demonstrate that throughput can be doubled or quadrupled without increasing power consumption or chip area. We survey Hopper's offerings related to these enhancements in \cref{subsec:hardware}. The principal technical hurdle is adapting \textsc{FlashAttention-2} so that it leverages these hardware advancements: asynchrony demands careful orchestration to overlap computations such as matmul and softmax despite their inherent dependencies, while employing low-precision requires judicious strategies to mitigate quantization errors, particularly when dealing with outlier activations in LLMs~\citep{dettmers2208llm, sun2024massive}.

For this purpose, we introduce \textsc{FlashAttention-3}, which integrates and advances three novel strategies aimed at enhancing efficiency on contemporary GPU architectures.\footnote{While our findings are presented with reference to NVIDIA’s Hopper architecture, our method remains applicable to any GPU hardware that offers strong asynchronous execution and supports low-precision operations.}

\begin{enumerate}

\item \textbf{Producer-Consumer asynchrony:} We propose a warp-specialized software pipelining method that leverages the asynchronous nature of both data transfer and Tensor Core computations by separating producer and consumer roles across distinct warps. This architectural separation enhances the algorithm’s capability to obscure both memory access and instruction issue delays.

\item \textbf{Hiding softmax under asynchronous block-wise GEMMs:} The low-throughput operations associated with softmax, such as floating-point multiply-add and exponentials, are scheduled to coincide with asynchronous WGMMA GEMM instructions, resulting in overlap between their execution. In implementing this, the \textsc{FlashAttention-2} algorithm is modified to avoid certain sequential constraints linking softmax and GEMM operations. As an illustration, in the algorithm’s 2-stage variant, while one block of the score matrix is being processed by softmax, an asynchronous WGMMA task computes the subsequent block in parallel.

\item \textbf{Hardware-accelerated low-precision GEMM:} The forward pass algorithm is refactored to facilitate GEMM on FP8 Tensor Cores, substantially boosting the observed TFLOPs/s, approximately doubling throughput. Achieving this necessitates reconciling layout conventions between WGMMA’s requirements for block storage of FP32 accumulators and FP8 operands. To address the decrease in accuracy caused by lowering precision to FP8, methods such as block quantization and incoherent processing are employed.
\end{enumerate}

To empirically assess our approach, we conducted extensive benchmarking of \textsc{FlashAttention-3} using the H100 SXM5 GPU across various configuration settings. The results demonstrate that: (1) FP16 delivers a 1.5–2.0$\times$ acceleration in the forward pass compared to \textsc{FlashAttention-2}, with throughput reaching up to 740 TFLOPs/s, and achieves 1.5–1.75$\times$ speedup in the backward pass; (2) FP8 attains nearly 1.2 PFLOPs/s; and (3) for extended sequence lengths, FP16 consistently surpasses competitors, while FP8 remains competitive\footnote{Specifically, for head dimension 64, FP8 \textsc{FlashAttention-3} leads, whereas for dimensions 128 and 256 it is comparable to the NVIDIA cuDNN attention for non-causal scenarios and trails slightly with causal masking.} relative to NVIDIA’s cuDNN library’s advanced attention implementation. Furthermore, FP16 \textsc{FlashAttention-3} matches \textsc{FlashAttention-2} in numerical error and outperforms conventional attention approaches, as intermediate operations—such as softmax rescaling—are maintained in FP32. Additionally, FP8 \textsc{FlashAttention-3}, when utilizing block quantization and incoherent computation strategies, achieves 2.6$\times$ higher accuracy than baseline attention using per-tensor quantization in situations with outlier features.

We have made \textsc{FlashAttention-3} publicly available under a permissive license\footnote{\textsc{FlashAttention-3}
  is available at \texttt{https://github.com/Dao-AILab/flash-attention}}, and we aim to incorporate its functionality into the PyTorch and Hugging Face ecosystems to maximize its accessibility and utility for researchers and developers worldwide.

\section{Background: Multi-Head Attention and GPU Characteristics}
\label{sec:background}

\subsection{Multi-Head Attention}
\label{subsec:multi_head_attn}

Let $\mathbf{Q}, \mathbf{K}, \mathbf{V} \in \mathbb{R}^{N \times d}$ be the query, key and value input sequences associated to a single head, where $N$ is the sequence length and $d$ is the head dimension. Then the attention output $\mathbf{O}$ is computed as:

\begin{equation*}
  \mathbf{S} = \alpha \mathbf{Q} \mathbf{K}^\top \in \mathbb{R}^{N \times N}, \quad \mathbf{P} = \mathrm{softmax}(\mathbf{S}) \in \mathbb{R}^{N \times N}, \quad \mathbf{O} = \mathbf{P}\mathbf{V} \in \mathbb{R}^{N \times d},
\end{equation*}

Here, $\mathrm{softmax}$ operates along the rows, and it is standard to use a scaling factor of $\alpha = 1/\sqrt{d}$. To mitigate issues of numerical instability caused by the exponential operation, $\mathrm{rowmax}(\mathbf{S})$ is subtracted from $\mathbf{S}$ prior to applying $\mathrm{softmax}$. In multi-head attention (MHA), distinct projection matrices for queries, keys, and values are allocated to each head. This process can be executed concurrently for several heads and across mini-batches, yielding the complete output tensor.

Now let $\phi$ be a scalar loss function and let $\mathbf{d}(-) = \partial \phi / \partial (-)$ be notation for the gradient. Given the output gradient $\mathbf{dO} \in \mathbb{R}^{N \times d}$, we compute $\mathbf{dQ}$, $\mathbf{dK}$, and $\mathbf{dV}$ according to the chain rule as follows:

\begin{align*}
  \mathbf{dV} &= \mathbf{P}^\top \mathbf{dO} \in \mathbb{R}^{N \times d} \\
  \mathbf{dP} &= \mathbf{dO} \mathbf{V}^\top \in \mathbb{R}^{N \times N} \\
  \mathbf{dS} &= \mathrm{dsoftmax} (\mathbf{dP}) \in \mathbb{R}^{N \times N} \\
  \mathbf{dQ} &= \alpha \mathbf{dS} \mathbf{K} \in \mathbb{R}^{N \times d} \\
  \mathbf{dK} &= \alpha \mathbf{dS}^\top \mathbf{Q} \in \mathbb{R}^{N \times d},
\end{align*}

In this context, $\mathbf{d}s = (\mathrm{diag}(p) - p p^\top)\mathbf{d}p$ where $p = \mathrm{softmax}(s)$, considering $s$ as a vector. The notation $\mathrm{dsoftmax}(\mathbf{dP})$ denotes applying this relation to each row individually. Ultimately, this process lends itself to parallelization across multiple heads and batches during the backward propagation phase of MHA.

\subsection{GPU hardware characteristics and execution model}
\label{subsec:hardware}

We outline elements of the GPU execution model pertinent to \textsc{FlashAttention-3}, emphasizing the NVIDIA Hopper architecture as a representative example of this framework.

\paragraph{Memory hierarchy:} The memory system in the GPU is structured as a tiered hierarchy of storage spaces, where higher levels provide less capacity but increased bandwidth (\cref{tab:gpu-hierarchy})\footnote{\citet{luo2024benchmarking} indicates a shared memory bandwidth of 128 bytes per clock cycle per SM; this figure is scaled by 132 SMs and a maximum clock speed of 1830 MHz for total bandwidth calculation.}. At the top, global memory (GMEM), sometimes referred to as HBM, consists of off-chip DRAM shared among all streaming multiprocessors (SMs). Data accessed from GMEM is automatically cached in an L2 cache on-chip. Moving further down, each SM features its own small, on-chip, programmer-managed cache called shared memory (SMEM), designed with high banking. The lowest level is comprised of the register file present within each SM.

\paragraph{Thread hierarchy:} Within the GPU programming paradigm, execution units are systematically arranged into a hierarchy of threads. Starting at the smallest granularity, individual threads are grouped into warps (each containing 32 threads), which are further combined into warpgroups (each consisting of 4 adjacent warps). These in turn form threadblocks (also known as cooperative thread arrays or CTAs), with higher-level organization including threadblock clusters (specific to Hopper), and ultimately, the grid.

These two hierarchical structures have a strong interdependence. Threads belonging to an identical CTA are simultaneously executed on a single SM, while CTAs grouped within the same cluster are concurrently mapped to the same GPC. The SMEM is accessible to every thread in a CTA, while each individual thread is restricted to accessing up to 256 registers (RMEM), which remain exclusive to that thread.

\paragraph{Asynchrony and warp-specialization:}

GPUs achieve high throughput by leveraging parallelism and asynchronous execution to mask both execution and memory access delays. To facilitate asynchronous data transfers between GMEM and SMEM, Hopper incorporates a specialized hardware component called the Tensor Memory Accelerator (TMA) \cite[\S7.29]{cuda}. In addition, the Hopper architecture advances beyond predecessors like Ampere by introducing an asynchronous Tensor Core, accessible through the warpgroup-scoped WGMMA instruction \cite[\S9.7.14]{ptx}, which is capable of reading inputs directly from shared memory.

The inclusion of hardware mechanisms for asynchrony facilitates the implementation of warp-specialized kernels, in which the warps within a CTA are assigned exclusively as either producers or consumers, responsible solely for data transfer or computational tasks, respectively. This approach broadly enhances the compiler’s capacity to produce highly efficient instruction schedules \citep{warp-specialization-2011}. Furthermore, Hopper offers dynamic register reallocation among warpgroups using \verb|setmaxnreg| \cite[\S9.7.17.1]{ptx}, enabling warps that execute MMAs to access a greater portion of RMEM, as opposed to warps dedicated to issuing TMA, which require only a single thread.

\paragraph{Low-precision number formats:}
\label{sec:low-precision-gpu}
Contemporary GPUs include dedicated hardware designed for efficient low-precision operations. As an illustration, the WGMMA instruction leverages FP8 Tensor Cores in the Hopper architecture, enabling each streaming multiprocessor (SM) to achieve double the performance relative to FP16 or BF16 computations.

To use FP8 WGMMA correctly, it is essential to be aware of operand layout requirements. For example, in a GEMM operation multiplying $A \times B^{\top}$, with $A$ being an $M\times K$ matrix and $B$ an $N\times K$ matrix, the $A$ or $B$ operand is deemed \emph{mn-major} if the memory layout is contiguous along the outer $M$ or $N$ axis, and \emph{k-major} when the contiguity is along the inner $K$ axis. FP16 WGMMA accommodates both mn-major and k-major configurations for operands residing in SMEM. In contrast, FP8 WGMMA exclusively supports operands with k-major alignment. This limitation complicates the use of dependent FP8 WGMMAs within a single kernel, particularly in cases such as attention, where consecutive GEMM operations must be fused and mismatches between FP32 accumulator layouts and FP8 operand arrangements create challenges.

Within the realm of attention mechanisms, such layout constraints necessitate specific adjustments to the development of an FP8 algorithm, a topic we elaborate on in \cref{sec:algofp8}.

\subsection{Standard Attention and Flash Attention} In accordance with \citet{dao2022flashattention}, we define \textbf{standard attention} as the traditional GPU-based approach in which the intermediate matrices $\mathbf{S}$ and $\mathbf{P}$ are explicitly stored in HBM. The principal innovation behind \textsc{FlashAttention} consists of employing a local form of the softmax reduction, which eliminates the need for costly memory operations associated with these intermediates, while merging the attention computation into a unified kernel. This notion of local softmax reduction is implemented in lines \ref{code-ws:softmax_start}-\ref{code-ws:softmax_end} of the consumer mainloop, as described in \cref{alg:flash3_wgmma_ws_only}, in conjunction with the scaling operations performed on blocks of $\mathbf{O}$. A concise derivation demonstrating that this approach indeed yields $\mathbf{O}$ is available in \cite[\S 2.3.1]{dao2023flashattention2}.

\section{FlashAttention-3: Algorithm}
\label{sec:algo}

This section introduces the \textsc{FlashAttention-3} algorithm. To clarify our exposition, we concentrate initially on the forward computation, reserving the discussion of the backward computation for~\cref{sec:algo_ws_bwd}. We begin by outlining how warp-specialization, combined with a circular SMEM buffer, can be integrated into the fundamental \textsc{FlashAttention-2} algorithm. Next, we delineate how to leverage the asynchronous capabilities of WGMMA to construct a 2-stage pipeline in which GEMM and softmax computations are overlapped. Lastly, we detail the changes required to support FP8, encompassing both layout compliance and accuracy improvements through block quantization and processing methods that are not mutually coherent.

\subsection{Producer-Consumer asynchrony through warp-specialization and pingpong scheduling}
\label{sec:algo_ws}

\paragraph{Warp-specialization}  
Similar to \textsc{FlashAttention-2}, the forward propagation in \textsc{FlashAttention-3} lends itself naturally to parallelization across the batch dimension, number of heads, and query sequence length. Therefore, we focus on describing the algorithm from a CTA perspective, where the computation targets a specific tile $\mathbf{Q}_i$ of the query matrix in order to produce the related output tile $\mathbf{O}_i$. For clarity, the warp-specialization approach described here makes use of a circular SMEM buffer but omits additional GEMM-softmax overlapping mechanisms. We denote by $d$ the dimension of each head and by $N$ the length of the sequence. Given a query block size $B_r$, we partition $\mathbf{Q}$ into $T_r = \lceil \frac{N}{B_r} \rceil$ segments, which are denoted $\mathbf{Q}_1, .., \mathbf{Q}_{T_r}$.

\begin{algorithm}[H]
    \caption{\small\label{alg:flash3_wgmma_ws_only}\textsc{FlashAttention-3} Forward Pass \textbf{Excluding} Intra-consumer Overlap -- CTA Perspective}
    \begin{algorithmic}[1]
\REQUIRE Given matrices $\mathbf{Q}_i \in \mathbb{R}^{B_r \times d}$ and $\mathbf{K}, \mathbf{V} \in \mathbb{R}^{N \times d}$ residing in HBM, along with key block size $B_c$ where $T_c = \lceil \frac{N}{B_c} \rceil$.
\STATE Instantiate a pipeline construct tasked with managing barrier synchronization for an $s$-stage circular SMEM buffer.
\IF {assigned to producer warpgroup}
\STATE Release a designated amount of registers as prearranged.
\STATE Initiate data transfer of $\mathbf{Q}_i$ from HBM to shared memory.
\STATE When data transfer of $\mathbf{Q}_i$ completes, signal this availability to consumer processes.
\FOR{$0 \le j < T_c$}
    \STATE Stall until the $(j\,\%\,s)$ buffer stage is marked as consumed.
    \STATE Begin loading blocks $\mathbf{K}_j, \mathbf{V}_j$ from HBM into the $(j\,\%\,s)$ buffer stage in shared memory.
    \STATE Upon successful loading, inform consumers that $\mathbf{K}_j, \mathbf{V}_j$ are now resident.
\ENDFOR
\ELSE \STATE Allocate required number of registers contingent on the number of consumer warps present.
\STATE Initialize the following locally in fast memory: $\mathbf{O}_i = (0) \in \mathbb{R}^{B_r \times d}$, $\ell_i = 0$ and $m_i = -\infty$ in $\mathbb{R}^{B_r}$.
\STATE Stall execution until $\mathbf{Q}_i$ is present in shared memory.
\FOR{$0 \le j < T_c$}
\STATE Stall until $\mathbf{K}_j$ transfer is completed into shared memory.
\STATE Compute $\mathbf{S}_i^{(j)} = \mathbf{Q}_i \mathbf{K}_j^T$ using SS-GEMM and commit result, blocking if necessary. \label{alg:ws_only_gemm1}
\STATE Save $m_i^{\mathrm{old}} = m_i$; update $m_i = \mathrm{max}(m_i^{\mathrm{old}}, \mathrm{rowmax}(\mathbf{S}_i^{(j)}))$. \label{code-ws:softmax_start}
\STATE Calculate $\widetilde{\mathbf{P}}_i^{(j)} = \mathrm{exp}(\mathbf{S}_i^{(j)} - m_i)$, and $\ell_i = \mathrm{exp}(m_i^{\mathrm{old}} - m_i) \ell_i + \mathrm{rowsum}(\widetilde{\mathbf{P}}_i^{(j)})$. \label{code-ws:softmax_end}
\STATE Wait for $\mathbf{V}_j$ to become ready in shared memory.
\STATE Update output: $\mathbf{O}_i = \mathrm{diag}(\mathrm{exp}(m_i^{\mathrm{old}} - m_i))^{-1} \mathbf{O}_i + \widetilde{\mathbf{P}}_i^{(j)} \mathbf{V}_j$ via RS-GEMM, then commit. \label{alg:ws_only_gemm2}
\STATE Notify producer by freeing the $(j\,\%\,s)$ buffer stage.
\ENDFOR
\STATE Normalize the output as $\mathbf{O}_i = \mathrm{diag}(\ell_i)^{-1} \mathbf{O}_i$; compute $L_i = m_i + \log(\ell_i)$.
\STATE Store results: copy $\mathbf{O}_i$ and $L_i$ back to HBM as block $i$ of $\mathbf{O}$ and $L$.
\ENDIF
\end{algorithmic}
\end{algorithm}

In our deployment of \cref{alg:flash3_wgmma_ws_only} on Hopper, we rely on \verb|setmaxnreg| for managing (de)allocations, TMA to handle data transfers for $\mathbf{Q}_i$ as well as $\{ \mathbf{K}_j, \mathbf{V}_j \}_{0 \leq j < T_c}$, and WGMMA for carrying out GEMMs within the consumer mainloop. The SS and RS prefixes specify if the first GEMM operand originates from shared memory or the register file, respectively. To properly follow the execution sequence in \cref{alg:flash3_wgmma_ws_only}, observe that TMA load operations proceed asynchronously, meaning the initiation of new loads does not depend on the completion of prior loads. Additionally, within the producer mainloop, buffer filling occurs during the first $s$ iterations, eliminating the need for issuing wait commands at these stages.

\paragraph{Pingpong scheduling} The asynchronous characteristics of WGMMA and TMA, combined with warp-specialization, facilitate concurrent execution—enabling the softmax calculations for one warpgroup to be performed simultaneously with the GEMM operations of another. This approach is particularly compelling when considering that, on contemporary hardware accelerators, non-matmul computations exhibit significantly reduced throughput relative to matmul computations. For instance, the H100 SXM5 GPU achieves 989 TFLOPS in FP16 matmul operations, yet only reaches 3.9 TFLOPS for special functions like exponentiation\footnote{According to the CUDA programming guide, each streaming multiprocessor (SM) can handle 16 special function operations per clock cycle. Using 132 SMs at a frequency of 1830 MHz yields 3.9 TFLOPS for these functions.}, which are required in the softmax procedure. During the attention forward computation in FP16 mode with a head dimension of 128, the total FLOPS for matmul surpasses those for exponentiation by a factor of 512. Nonetheless, the throughput for exponentiation is 256 times less, implying that exponentiation may occupy 50\% of the computation cycle allotted to matmul. The disparity intensifies when working with FP8, since the throughput of matmul doubles while exponentiation throughput remains static.

The computation of the exponential function is handled by a dedicated hardware component known as the multi-function unit, and therefore, it is optimal to overlap this operation with matrix multiplication tasks executed by the Tensor Cores. To orchestrate this concurrency, we utilize synchronization barriers (\texttt{bar.sync} instructions) to ensure that the GEMM operations (specifically, GEMM1 -- $\mathbf{P} \mathbf{V}$ within one iteration, and GEMM0 -- $\mathbf{Q} \mathbf{K}^\top$ from the following iteration) executed by warpgroup 1 are completed prior to those of warpgroup 2. This arrangement enables the softmax computation associated with warpgroup 1 to proceed concurrently as warpgroup 2 performs its GEMMs. Subsequently, the warpgroup responsibilities are interchanged: warpgroup 2 executes softmax while warpgroup 1 undertakes its GEMM computations, resulting in a so-called “pingpong” execution pattern. This mechanism is visualized in~\cref{fig:pingpong_scheduling}. Although the actual scheduling deviates somewhat from the schematic representation, we typically observe notable performance benefits (for instance, increasing throughput from 570 TFLOPS to 620-640 TFLOPS for FP16 forward computations with a head dimension of 128 and sequence length of 8192).

\paragraph{Attention variants} In the case of multi-query attention~\citep{shazeer2019fast} and grouped query attention~\citep{ainslie2023gqa}, we adopt the strategy used by \textsc{FlashAttention-2}, modifying tensor indexing to prevent redundant storage of $\mathbf{K}$ and $\mathbf{V}$ within HBM.

\subsection{Intra-warpgroup overlapping GEMMs and softmax}

Within a single warpgroup, it is possible to execute certain instructions from the softmax operation concurrently with those from the GEMMs. Here, we outline a method to achieve this overlap.

In the attention algorithm, operations within the inner loop (main loop) have sequential dependencies that impede parallelization within a single iteration. For example, (local) softmax (lines \ref{code-ws:softmax_start} to \ref{code-ws:softmax_end}) relies on the output $\mathbf{S}_i^{(j)}$ of the first GEMM, while the second GEMM takes its result $\widetilde{\mathbf{P}}_i^{(j)}$ as an operand. Indeed, the wait statements in lines \ref{alg:ws_only_gemm1} and \ref{alg:ws_only_gemm2} of \cref{alg:flash3_wgmma_ws_only} serialize the execution of softmax and GEMMs. However, we can break these dependencies by pipelining across iterations through additional buffers in registers. Pursuing this idea, we propose the following two-stage\footnote{Note that the number of stages of the overlapping scheme is bounded by, but need not equal, the number $s$ of stages in the circular SMEM buffer.} GEMM-softmax pipelining algorithm:

\begin{algorithm}[H]
    \caption{\small\label{alg:flash3_wgmma}\textsc{FlashAttention-3} consumer warpgroup forward pass}
    \begin{algorithmic}[1]
      \REQUIRE  Matrices $\mathbf{Q}_i \in \mathbb{R}^{B_r \times d}$ and $\mathbf{K}, \mathbf{V} \in \mathbb{R}^{N \times d}$ stored in HBM, along with key block dimension $B_c$, and $T_c = \lceil \frac{N}{B_c} \rceil$.
      \STATE Assign a pre-calculated number of registers as determined by the number of consumer warps.
      \STATE Allocate $\mathbf{O}_i = (0) \in \mathbb{R}^{B_r \times d}$, and vectors $\ell_i, m_i = (0), (-\infty) \in \mathbb{R}^{B_r}$ in on-chip memory.
      \STATE Await the transfer of $\mathbf{Q}_i$ and initial block $\mathbf{K}_0$ into shared memory.
      \STATE Perform the operation $\mathbf{S}_{\mathrm{cur}} = \mathbf{Q}_i \mathbf{K}_0^T$ via WGMMA. After submitting, wait for completion.
      \STATE Free the buffer's stage $0$ corresponding to $\mathbf{K}$.
      \STATE Using $\mathbf{S}_{\mathrm{cur}}$, update $m_{i}$, compute $\tilde{\mathbf{P}}_{\mathrm{cur}}$, and $\ell_{i}$, followed by rescaling $\mathbf{O}_i$.
      \FOR{$1 \le j < T_c - 1$}
        \STATE Await the transfer of block $\mathbf{K}_j$ into shared memory.
        \label{code:mainloop_start}
        \STATE Calculate $\mathbf{S}_{\mathrm{next}} = \mathbf{Q}_i \mathbf{K}_{j}^T$ using WGMMA. Submit without waiting. \label{code:first_wgmma}
        \STATE Wait until $\mathbf{V}_{j-1}$ has been loaded in shared memory.
        \STATE Update $\mathbf{O}_{i} = \mathbf{O}_{i} + \tilde{\mathbf{P}}_{\mathrm{cur}} \mathbf{V}_{j-1}$ employing WGMMA. Submit without waiting. \label{code:second_wgmma}
        \STATE Wait for the result of WGMMA operation $\mathbf{Q}_i \mathbf{K}_{j}^T$.
        \STATE Using $\mathbf{S}_{\mathrm{next}}$, calculate $m_{i}$, $\tilde{\mathbf{P}}_{\mathrm{next}}$, and $\ell_{i}$. \label{code:softmax}
        \STATE Wait for the WGMMA result for $\tilde{\mathbf{P}}_{\mathrm{cur}} \mathbf{V}_{j-1}$; rescale $\mathbf{O}_i$ afterward.
        \STATE Release the $(j\,\%\,s)$th and $(j-1\,\%\,s)$th buffer stages for $\mathbf{K}$ and $\mathbf{V}$, respectively.
        \STATE Assign $\mathbf{S}_{\mathrm{cur}}$ to be the contents of $\mathbf{S}_{\mathrm{next}}$.
        \label{code:mainloop_end}
      \ENDFOR \STATE Await loading of $\mathbf{V}_{T_c - 1}$ in shared memory.
      \STATE Calculate 
      $\mathbf{O}_{i} = \mathbf{O}_{i} + \tilde{\mathbf{P}}_{\mathrm{last}} \mathbf{V}_{T_c - 1}$ with WGMMA. Commit, then wait for result.

\STATE Epilogue: Adjust $\mathbf{O}_{i}$ according to $m_{i}$. Determine $L_{i}$ from $m_{i}$ and $\ell_{i}$.  
Store $\mathbf{O}_{i}$ and $L_{i}$ in HBM, assigning them to the $i$-th segments of $\mathbf{O}$ and $L$.  
\end{algorithmic}
\end{algorithm}

\cref{alg:flash3_wgmma} serves as a substitute for the consumer segment in \cref{alg:flash3_wgmma_ws_only}, thus forming the entirety of the \textsc{FlashAttention-3} algorithm operating at FP16 precision. Conceptually, WGMMA is utilized as shorthand for asynchronous GEMM throughout the discussion. In the main processing loop (spanning lines \ref{code:mainloop_start} to \ref{code:mainloop_end}), the second WGMMA invocation in iteration $j$ (specified in line \ref{code:second_wgmma}) is executed concurrently with the softmax computations from iteration $j+1$ (located at line \ref{code:softmax}).

Although the pipelined architecture depicted above suggests potential theoretical performance improvements, several practical challenges arise in real-world scenarios:

\paragraph{Compiler reordering} While the pseudocode conveys an ideal sequence of execution, the NVCC compiler frequently reorders instructions to achieve better optimization. Such reordering can interfere with the intended interleaving of WGMMA and non-WGMMA operations, possibly resulting in altered behavior or reduced expected speedup. Nevertheless, an inspection of the generated SASS code indicates that the compiler maintains the required overlapping, as further discussed in Section~\ref{sec:2-stage-sass}.

\paragraph{Register pressure} Achieving optimal throughput demands that register spilling be kept to a minimum. Yet, the two-stage pipeline introduces extra pressure on registers, as it needs to allocate additional resources for intermediary storage and stage-specific context. In particular, the pipeline must reserve space for an additional $\mathbf{S}_{\mathrm{next}}$ variable in registers, accounting for an increase of $B_r \times B_c \times \text{sizeof}(\text{float})$ of register usage per threadblock. This heightened usage might clash with strategies aiming to increase the block size—another commonly pursued optimization—which also intensifies register requirements. Consequently, such conflicts necessitate careful performance profiling to decide on appropriate trade-offs.

\paragraph{3-stage pipelining} As an extension to the aforementioned 2-stage method, we introduce a 3-stage pipeline variant designed to additionally overlap a second WGMMA computation with the softmax operation. This technique holds promise for further increasing Tensor Core utilization. However, introducing a third stage also heightens the register requirements, intensifying the difficulty in managing the trade-off between the depth of pipelining and the tile size. An expanded explanation of this 3-stage approach, along with an assessment of its impact, is provided in ~\cref{sec:3-stage}.

\subsection{Low-precision with FP8}
\label{sec:algofp8}

\textbf{Efficiency: layout transformations.} Performing the forward computation of \textsc{FlashAttention-3} in FP8 precision introduces unique difficulties related to layout compatibility, which are not present when using FP16.

To begin, it is important to recognize that the input tensors $\mathbf{Q}$, $\mathbf{K}$, and $\mathbf{V}$ are conventionally laid out with contiguous elements in the head dimension. However, adhering to the k-major ordering required by FP8 WGMMA for the second GEMM necessitates that $\mathbf{V}$—or more precisely, the tiles of $\mathbf{V}$ that are placed into SMEM—should be contiguous along the sequence length dimension. Given that the TMA load operation preserves the contiguity of dimensions and cannot modify it, this presents two possible approaches: (1) performing a transpose of $\mathbf{V}$ within GMEM as a preprocessing stage, or (2) executing an in-kernel transpose of each tile of $\mathbf{V}$ once they have been moved into SMEM. To realize the first strategy, one can either (1a) incorporate the transpose into the epilogue of an earlier operation such as rotary embedding, or (1b) invoke an independent preprocessing transpose kernel\footnote{A well-tuned transpose kernel is capable of matching the device's bandwidth performance~\citep{colfax_cutlass_transpose_2024}.} to swap the strides between the sequence length and head dimensions. Yet, integrating approach (1a) with standard libraries proves to be challenging, while option (1b) is not optimal in memory-constrained contexts like inference due to its excessive resource consumption.

For FP8 \textsc{FlashAttention-3}, we adopt approach (2). Within the kernel, the transpose operation leverages the LDSM (\verb|ldmatrix|) and STSM (\verb|stmatrix|) instructions, which permit a warp to cooperatively transfer data between shared memory (SMEM) and register memory (RMEM) in blocks of 128 bytes. \footnote{The PTX documentation specifies LDSM/STSM as handling $8 \times 8$ matrices of 16-bit values \cite[\S 9.7.13.4.15-16]{ptx}, but we achieve FP8 usage by packing pairs of 8-bit values together for these instructions. Nonetheless, since LDSM/STSM transpose variants are unable to decompose packed 8-bit values, intermediate register shuffling is required between LDSM and STSM to accomplish proper tile-wise transposition; we omit the specifics here.} Because LDSM/STSM are register-friendly, they can be efficiently run within the producer warpgroup while also enabling layout transposition during memory movement. Additionally, once the initial pass is done, we can schedule the transposition of the upcoming $\mathbf{V}$ tile to overlap with the shadow computation of two WGMMAs that process the preceding $\mathbf{V}$ and current $\mathbf{K}$ tiles.

Second, we note a distinction between the memory layout used for the FP32 accumulator in an FP8 WGMMA and the one assumed for operand A when stored in registers, unlike the behavior observed with FP16. Figures \cref{fig:rmem_accum} and \cref{fig:rmem_operand} illustrate fragments of these layouts, showing how register entries are distributed per thread in the specified order. By utilizing byte permute instructions, the accumulator resulting from the first WGMMA can be reorganized into a layout appropriate for the subsequent WGMMA, which also aligns with the configuration of the $\mathbf{V}$ tile generated by the in-kernel transpose. Referring to \cref{fig:rmem_accum}, we apply a sequence transformation such that
$$\{ \verb|d0 d1 d4 d5 d2 d3 d6 d7| \},$$
with this register arrangement applied repeatedly for every block of 8 bytes. From the perspective of the logical configuration of the $\mathbf{P}$ tile, this operation reorders its columns—for instance, columns $0189$ are shifted to the first four positions. To enable WGMMA to produce the intended output tile, the in-kernel transpose can correspondingly output a row permutation for the $\mathbf{V}$ tile.\footnote{This flexibility in performing the transpose inside the kernel obviates the necessity for shuffle instructions that would change register allocation across threads, as previously detailed in~\citep{colfax_fp8_flashattention_2024}.}

\textbf{Accuracy: block quantization and incoherent processing.} The FP8 (e4m3) numerical representation encodes only 3 bits for the mantissa and 4 bits for the exponent. This limited precision induces greater numerical discrepancies compared to FP16 or BF16 formats. Additionally, the distribution of values in large-scale models often exhibits extreme outliers~\citep{dettmers2208llm,
  sun2024massive}, which pose challenges for effective quantization due to their disproportionately large magnitude. A common approach involves per-tensor scaling~\citep{micikevicius2022fp8}, assigning a distinct scaling factor to each tensor (such as for $\mathbf{Q}$, $\mathbf{K}$, or $\mathbf{V}$).
To mitigate numerical inaccuracies arising from FP8 quantization during attention computation, we leverage two strategies:

\begin{enumerate}

\item \textbf{Block quantization}: This method involves retaining a single scalar value for each block, where for the matrices $\mathbf{Q}$, $\mathbf{K}$, and $\mathbf{V}$, we partition the tensors into blocks with dimensions $B_r \times d$ or $B_c \times d$ and quantize the blocks independently. This quantization step can be seamlessly integrated with preprocessing operations (such as rotary embedding) immediately preceding the attention mechanism, without incurring any latency penalty due to memory bandwidth constraints of rotary embedding. Additionally, as the \textsc{FlashAttention-3} algorithm inherently processes data in block-wise fashion, each block of $\mathbf{S}$ may be rescaled to offset block quantization effects without imposing computational overhead.

\item \textbf{Incoherent processing}: To mitigate the influence of outlier values, we pre-process both $\mathbf{Q}$ and $\mathbf{K}$ by multiplying them with a random orthogonal matrix $\mathbf{M}$ prior to FP8 quantization. Because $\mathbf{M}$ is orthogonal and satisfies $\mathbf{M} \mathbf{M}^\top = I$, the transformation maintains the result of the attention computation, i.e., $(\mathbf{Q} \mathbf{M}) (\mathbf{K} \mathbf{M})^\top = \mathbf{Q} \mathbf{K}^\top$. This randomization redistributes outlier effects, since each element in $\mathbf{Q}\mathbf{M}$ or $\mathbf{K}\mathbf{M}$ is formed as a random linear combination of the original matrix elements, thereby attenuating quantization inaccuracies. Following the practices in \citet{chee2024quip} and~\citet{tseng2024quip}, we select $\mathbf{M}$ as a product of randomly chosen diagonal matrices with entries $\pm 1$ and a Hadamard matrix, enabling multiplication in $O(d \log d)$ time as opposed to $O(d^2)$, which may also be combined with rotary embedding at no extra computational cost.

\end{enumerate}

Empirical results, presented in \cref{sec:numerical_error}, demonstrate that these approaches can decrease numerical error by up to a factor of 2.6.

\section{Empirical Validation}
\label{sec:exp}

Our implementation of \textsc{FlashAttention-3} leverages CUTLASS~\citep{Thakkar_CUTLASS_2023} primitives, notably the WGMMA and TMA abstractions, for assessing both the performance and precision of the algorithm.

\begin{itemize}

\item \textbf{Attention Performance Evaluation.} We assess the execution time of \textsc{FlashAttention-3} at varying sequence lengths, juxtaposing its performance with the conventional PyTorch implementation, \textsc{FlashAttention-2}, the Triton-based \textsc{FlashAttention-2} (which leverages H100-specialized instructions), and a vendor-provided cuDNN version of \textsc{FlashAttention-2} tailored for H100 GPUs. Our analysis demonstrates that \textsc{FlashAttention-3} achieves speed gains of up to 2.0$\times$ over \textsc{FlashAttention-2}, and 1.5$\times$ over the Triton variant. Additionally, \textsc{FlashAttention-3} attains throughput rates up to 740 TFLOPs/s, amounting to 75\% of the H100 GPU's peak theoretical TFLOPs/s.
\item \textbf{Ablation Experiments.} We further show that enhancements such as warp-specialization and GEMM-softmax pipelining significantly impact the efficiency gains seen in \textsc{FlashAttention-3}.
\item \textbf{FP8 Attention Precision.} We demonstrate that employing block quantization along with incoherent processing lowers the numerical inaccuracies in FP8 \textsc{FlashAttention-3} by a factor of 2.6$\times$.
\end{itemize}

\subsection{Benchmarking Attention}
\label{subsec:benchmark_attn}

We evaluate the execution time of several attention mechanisms utilizing an H100 80GB SXM5 GPU, under varying configurations—such as with and without causal masking and using head dimensions of either 64 or 128—on FP16 input data. The outcomes, presented in~\cref{fig:benchmark_attn_fwd} and~\cref{fig:benchmark_attn_bwd}, reveal that \textsc{FlashAttention-3} achieves approximately 1.5 to 2.0$\times$ speedup over \textsc{FlashAttention-2} for the forward pass, and attains 1.5 to 1.75$\times$ improvement during the backward pass. In comparison to conventional attention implementations, \textsc{FlashAttention-3} is capable of delivering a speed increase ranging from 3 to 16$\times$. Furthermore, for sequence lengths of 1,000 tokens or more, \textsc{FlashAttention-3} even outperforms the proprietary, GPU-specialized library cuDNN.

\paragraph{Benchmark settings:} We vary the sequence length as 512, 1k, ..., 16k, and set batch size so that the total number of tokens is 16k. We set the hidden dimension to 2048, and head dimension to be either 64, 128, or 256 (i.e., 32 heads, 16 heads, or 8 heads). To calculate the FLOPs of the forward pass, we use:

\begin{equation*}
  4 \cdot \text{seqlen}^2 \cdot \text{head dimension} \cdot \text{number of heads}.
\end{equation*}

By applying causal masking, we reduce the quantity by a factor of 2, reflecting that roughly half of the elements are evaluated. For estimating the FLOPs during the backward pass, the forward pass FLOPs are scaled by 2.5, as the backward pass involves five matrix multiplications compared to the two performed in the forward pass, which is attributable to recomputation.

Additionally, we evaluate the runtime of FP8 during the forward pass using comparable conditions. The findings for headdim 256 are displayed in ~\cref{fig:benchmark_attn_fp8}, with comprehensive results provided in \cref{sec:benchmark_attn_fp8_full}.

\subsection{Ablation Study: 2-Stage Pipelining Experiments}
\label{sec:2-stage-pipelining-experiments}

We conduct an ablation study on the 2-stage WGMMA-softmax pipelining and warp-specialization components for non-causal FP16 \textsc{FlashAttention-3}, maintaining fixed configuration $\{\text{batch}, \text{seqlen}, \text{nheads}, \text{hdim}\} = \{ 4, 8448, 16, 128\}$. Results presented in~\cref{table:ablation_pipelining} demonstrate that our algorithmic enhancements—specifically, introducing asynchrony via warp-specialization and achieving overlap between GEMM and softmax operations—yield a considerable performance improvement, raising throughput from 570 to 661 TFLOPs.

\subsection{Numerical Error Validation}
\label{sec:numerical_error}

Given growing concern regarding the numerical inaccuracies in \textsc{FlashAttention}~\citep{golden2024flash}, we carry out a comparison of \textsc{FlashAttention-2}, \textsc{FlashAttention-3}, and a conventional attention mechanism with respect to a high-precision FP64 reference implementation. To model the presence of extreme features and activations that can occur in large language models~\citep{dettmers2208llm, sun2024massive}, we initialize the values of $\mathbf{Q}, \mathbf{K}, \mathbf{V}$ according to the following probabilistic distribution:

\begin{equation*}
  \mathcal{N}(0, 1) + \mathcal{N}(0, 100) \cdot \mathrm{Bernoulli}(0.001).
\end{equation*}

In this setup, every element follows a normal distribution with a mean of zero and a standard deviation of 1, except for 0.1\% of the elements, where an additional independent Gaussian variable with a standard deviation of 10 is added. The root mean squared error (RMSE) is reported in~\cref{table:numerical_error}. When using FP16, both \textsc{FlashAttention-2} and \textsc{FlashAttention-3} attain an RMSE that is 1.7$\times$ lower than the conventional approach, as they maintain intermediate softmax computations in FP32. The FP8 baseline implementation applies per-tensor scaling, stores the matmul accumulator in FP32, and retains intermediate softmax outcomes in FP16. Leveraging block quantization and non-coherent processing, \textsc{FlashAttention-3} in FP8 demonstrates a 2.6$\times$ improvement in accuracy over this baseline.

\section{Dicussion, Limitations, Conclusion}
\label{sec:discussion}

Through the development of \textsc{FlashAttention-3}, we illustrate that innovative programming strategies and emerging hardware capabilities—including asynchrony and low-precision computation—significantly enhance both the speed and fidelity of attention mechanisms. Our approach achieves a 1.5-2.0$\times$ acceleration over \textsc{FlashAttention-2}, while also decreasing FP8 numerical errors by a factor of 2.6 relative to conventional per-tensor quantization methods. Some areas for future improvement include: tailoring the approach for LLM inference, incorporating a persistent kernel architecture into the FP8 kernel,\footnote{In our benchmarks, FP16 \textsc{FlashAttention-3} employs a persistent kernel and load balancing, unlike FP8 \textsc{FlashAttention-3}, which lacks these optimizations. This is a partial reason for FP8 \textsc{FlashAttention-3}'s suboptimal performance on smaller sequence lengths and causal masking tasks compared to FP8 cuDNN kernels.} and further exploring how low-precision attention influences training at scale. While our experimentation is based on Hopper GPUs, we anticipate that the principles and methods introduced can be generalized to other hardware accelerator platforms. Ultimately, we envisage that making attention both faster and more accurate will enable progress in applications that require extensive context handling.

\end{document}